\begin{abstract}
Đồ án này trình bày thiết kế, phân tích và triển khai một hệ thống khử nhiễu tiếng nói (Speech Enhancement) toàn dải (48 kHz) theo thời gian thực, dựa trên kiến trúc nơ-ron hồi quy kết hợp xử lý tín hiệu số (Hybrid DSP/Deep Learning) lấy cảm hứng từ mô hình RNNoise. Mục tiêu trọng tâm là hiện thực hóa thuật toán trên vi điều khiển nhúng ESP32-S3, đồng thời đánh giá chi tiết các rào cản vi kiến trúc (micro-architecture) của phần cứng. Thông qua việc phân tích phổ và mô phỏng thính giác Bark, mô hình AI dự đoán các hệ số khuếch đại (gain) lý tưởng, kết hợp cùng bộ lọc lược (comb filter) để bảo toàn cấu trúc hài âm. Bằng các kỹ thuật tối ưu hóa mức hệ thống như sử dụng tập lệnh SIMD PIE (Processor Instruction Extensions), định tuyến bộ nhớ IRAM cho các nhân tính toán và lượng tử hóa, độ trễ cấu phần DSP đã giảm đáng kể (lên đến 3.64 lần). Dù vậy, nghiên cứu chỉ ra rằng băng thông bộ nhớ PSRAM và hiện tượng Cache Thrashing trên bộ đệm L1 (32KB) cấu thành nút thắt cổ chai vật lý không thể vượt qua đối với mạng suy luận kích thước 1.3MB khi yêu cầu độ trễ dưới 10ms. Để khắc phục, một kiến trúc lai phân tán (Edge-to-Server) được đề xuất, trong đó ESP32-S3 đóng vai trò \textit{Audio Gateway} đa nhân, chuyển giao luồng dữ liệu cho máy chủ biên dịch bằng AVX2, giúp giảm tổng độ trễ suy luận xuống dưới 0.1ms và duy trì chất lượng thoại xuất sắc.
\end{abstract}
